\NormalFont\ShortTitle{Jonas}
{\MT Jonas

\par }\ChapOne{1}{\PP \VerseOne{1}E veio a palavra do SENHOR a Jonas, filho de Amitai, dizendo:
\VS{2}Levanta-te, vai à grande cidade de Nínive, e prega contra ela; porque sua maldade subiu diante de mim.
\VS{3}Porém Jonas se levantou para fugir da presença do SENHOR a Társis, e desceu a Jope; e achou um navio que partia para Társis. Então pagou sua passagem e entrou nele, para ir com eles a Társis
{\ADD{a fim de se distanciar}} do SENHOR.
\VS{4}Mas o SENHOR fez levantar um grande vento no mar; e fez-se uma tempestade tão grande no mar, que o navio estava a ponto de se quebrar.
\VS{5}Então os marinheiros tiveram medo, e cada um clamava a seu deus; e lançaram no mar os objetos que havia no navio, para com eles diminuir o peso. Jonas, porém, havia descido ao interior do navio, e se pôs a dormir profundamente.
\VS{6}E o capitão do navio se aproximou dele, e lhe disse: Que há contigo, dorminhoco? Levanta-te, e clama a teu deus; talvez ele se lembre de nós, para que não pereçamos.
\VS{7}E disseram cada um a seu companheiro: Vinde, e lancemos sortes, para sabermos por causa de quem este mal
{\ADD{veio}} sobre nós. Então lançaram sortes, e a sorte caiu sobre Jonas.
\VS{8}Então eles lhe disseram: Conta-nos, por favor, por que nos este mal nos
{\ADD{veio}} . Qual é a tua profissão, e de onde vens? Qual é a tua terra, e de que povo és?
\VS{9}E ele lhes respondeu: Sou hebreu, e temo ao SENHOR, o Deus dos céu, que fez o mar e a terra.
\VS{10}Então aqueles homens tiveram um grande medo, e lhe disseram: Por que fizeste isto? Pois aqueles homens sabiam que ele estava fugindo do SENHOR, porque ele já havia lhes dito.
\VS{11}Então lhe disseram: O que faremos contigo, para que o mar se nos aquiete?Porque o mar cada vez mais se embravecia.
\VS{12}E ele lhes respondeu: Tomai-me, e lançai-me ao mar, e o mar se vos aquietará; porque eu sei que foi por minha causa que esta grande tempestade
{\ADD{veio}} sobre vós.
\VS{13}Mas os homens se esforçavam para remar, tentando voltar à terra firme; porém não conseguiam, porque o mar cada vez mais se embravecia sobre eles.
\VS{14}Então clamaram ao SENHOR, e disseram: Ó SENHOR, nós te pedimos para não perecermos por causa da alma deste homem, nem ponhas sangue inocente sobre nós; porque, tu, SENHOR, fizeste como te agradou.
\VS{15}Então tomaram a Jonas, e o lançaram ao mar; e o mar se aquietou de seu furor.
\VS{16}Por isso aqueles homens temeram ao SENHOR
{\ADD{com}} grande temor; e ofereceram sacrifício ao SENHOR, e prometeram votos.
\VS{17}Mas o SENHOR havia preparado um grande peixe para que tragasse a Jonas; e Jonas esteve três dias e três noites no ventre do peixe.

\par }\Chap{2}{\PP \VerseOne{1}E Jonas orou ao SENHOR seu Deus, desde as entranhas do peixe,
\VS{2}E disse: Clamei da minha angústia ao SENHOR, e ele me respondeu; do ventre do Xeol
\FTNT{*}{{\FR 2:2 }
Xeol é o lugar dos mortos
} gritei,
{\ADD{e}} tu ouviste minha voz.
\VS{3}Pois tu me lançaste no profundo, no meio dos mares, e a correnteza me cercou; todas as tuas ondas e tuas vagas passaram sobre mim.
\VS{4}E eu disse: Lançado estou de diante de teus olhos; porém voltarei a ver o teu santo templo.
\FTNT{†}{{\FR 2:4 }
porém voltarei … templo
trad. alt. como voltarei a ver o teu santo templo?
}
\VS{5}As águas me cercaram até a alma, o abismo me rodeou; as algas se enrolaram à minha cabeça.
\VS{6}Desci aos fundamentos dos montes; a terra trancou suas fechaduras a mim para sempre; porém tu tiraste minha vida da destruição, ó SENHOR meu Deus.
\VS{7}Quando minha alma desfalecia em mim, eu me lembrei do SENHOR; e minha oração chegou a ti em teu santo templo.
\VS{8}Os que dão atenção a coisas inúteis e ilusórias abandonam sua própria misericórdia.
\VS{9}Mas eu porém sacrificarei a ti com voz de gratidão; pagarei o que prometi. A salvação
{\ADD{vem}} do SENHOR.
\VS{10}E o SENHOR falou ao peixe, e este vomitou a Jonas em terra firme.

\par }\Chap{3}{\PP \VerseOne{1}E a palavra do SENHOR veio pela segunda vez a Jonas, dizendo:
\VS{2}Levanta-te, e vai à grande cidade de Nínive, e prega contra ela a pregação que eu te falo.
\VS{3}Então Jonas se levantou, e foi a Nínive, conforme a palavra do SENHOR. E Nínive era cidade extremamente grande, de três dias de caminho.
\VS{4}E Jonas começou a entrar pela cidade, caminho de um dia, e pregava dizendo: Daqui a quarenta dias Nínive será destruída.
\VS{5}E os homens de Nínive creram em Deus; e convocaram um jejum, e se vestiram de sacos, desde o maior até o menor deles.
\VS{6}Pois quando esta palavra chegou ao rei de Nínive, ele se levantou de seu trono, e lançou de si sua roupa, e cobriu-se de saco, e se sentou em cinzas.
\VS{7}E convocou e anunciou em Nínive, por mandado do rei e de seus maiorais, dizendo: Nem pessoas nem animais, nem bois nem ovelhas, provem coisa alguma, nem se lhes dê alimento, nem bebam água;
\VS{8}E que as pessoas e os animais se cubram de saco, e clamem fortemente a Deus; e convertam- se cada um de seu mau caminho, e da violência que está em suas mãos.
\VS{9}Quem sabe Deus mude de ideia e se arrependa, e se afaste do ardor de sua ira, para não perecermos?
\VS{10}E Deus viu a atitude deles, que se converteram de seu mau caminho; então Deus se arrependeu do mal que tinha dito que lhes faria, e não o fez.

\par }\Chap{4}{\PP \VerseOne{1}Mas Jonas se desgostou muito, e se encheu de ira.
\VS{2}E orou ao SENHOR, e disse: Ah, SENHOR, não foi isto o que eu dizia enquanto ainda estava em minha terra? Por isso me preveni fugindo a Társis; porque sabia eu que tu és Deus gracioso e misericordioso, que demoras a te irar, tens grande misericórdia, e te arrependes do mal.
\VS{3}Agora pois, ó SENHOR, peço-te que me mates, porque é melhor para mim morrer do que viver.
\VS{4}E o SENHOR lhe disse: É correta
\FTNT{*}{{\FR 4:4 }
correta
trad. alt. tanta
} essa tua ira?
\VS{5}E Jonas saiu da cidade, e se sentou ao oriente da cidade, e fez para si ali uma barraca, e se sentou debaixo dela à sombra, até ver o que seria da cidade.
\VS{6}E o SENHOR Deus preparou uma planta, e a fez crescer sobre Jonas para que fizesse sombra sobre sua cabeça, e lhe aliviasse de seu mal-estar; e Jonas se alegrou grandemente por causa da planta.
\VS{7}Mas Deus preparou uma lagarta no dia seguinte ao amanhecer, a qual feriu a planta, e ela se secou.
\VS{8}E aconteceu que, ao levantar do sol, Deus preparou um quente vento oriental; e o sol feriu a Jonas na cabeça, de modo que ele se desmaiava, e desejava a morte, dizendo: Melhor é para mim morrer do que viver.
\VS{9}Então Deus disse a Jonas: É correta
\FTNT{†}{{\FR 4:9 }
correta
trad. alt. tanta
} a tua ira pela planta?E ele respondeu: É correto eu me irar
\FTNT{‡}{{\FR 4:9 }
É correto eu me irar
trad. alt. Estou muito irado
} até a morte.
\VS{10}E o SENHOR disse: Tu tiveste pena da planta, na qual não trabalhaste, nem tu a fizeste crescer, que em uma noite nasceu, e em outra noite pereceu;
\VS{11}E não teria eu pena de Nínive, aquela grande cidade onde há mais de cem e vinte mil pessoas que não sabem a diferença entre sua mão direita e a esquerda, e muitos animais?\par }